\documentclass[10pt,a4paper]{amsart}
\usepackage{amsmath,amssymb,amsthm,graphicx}
% For commutative diagrams we use Xy-pic
\usepackage[all]{xy}

% We use 2cell for 2-commutative diagrams.
\xyoption{2cell}
% Theorem environments.
%
\theoremstyle{plain}
\newtheorem{theorem}[subsection]{Theorem}
\newtheorem{proposition}[subsection]{Proposition}
\newtheorem{lemma}[subsection]{Lemma}

\theoremstyle{definition}
\newtheorem{definition}[subsection]{Definition}
\newtheorem{example}[subsection]{Example}
\newtheorem{exercise}[subsection]{Exercise}
\newtheorem{situation}[subsection]{Situation}

\theoremstyle{remark}
\newtheorem{remark}[subsection]{Remark}
\newtheorem{remarks}[subsection]{Remarks}

\numberwithin{equation}{subsection}

% Macros
%
\def\lim{\mathop{\mathrm{lim}}\nolimits}
\def\colim{\mathop{\mathrm{colim}}\nolimits}
\def\Spec{\mathop{\mathrm{Spec}}}
\def\Hom{\mathop{\mathrm{Hom}}\nolimits}
\def\Ext{\mathop{\mathrm{Ext}}\nolimits}
\def\SheafHom{\mathop{\mathcal{H}\!\mathit{om}}\nolimits}
\def\SheafExt{\mathop{\mathcal{E}\!\mathit{xt}}\nolimits}
\def\Sch{\mathit{Sch}}
\def\Mor{\mathop{\mathrm{Mor}}\nolimits}
\def\Ob{\mathop{\mathrm{Ob}}\nolimits}
\def\Sh{\mathop{\mathit{Sh}}\nolimits}
\def\NL{\mathop{N\!L}\nolimits}
\def\CH{\mathop{\mathrm{CH}}\nolimits}
\def\proetale{{pro\text{-}\acute{e}tale}}
\def\etale{{\acute{e}tale}}
\def\QCoh{\mathit{QCoh}}
\def\Ker{\mathop{\mathrm{Ker}}}
\def\Im{\mathop{\mathrm{Im}}}
\def\Coker{\mathop{\mathrm{Coker}}}
\def\Coim{\mathop{\mathrm{Coim}}}

% Boxtimes
%
\DeclareMathSymbol{\boxtimes}{\mathbin}{AMSa}{"02}

%
% Macros for moduli stacks/spaces
%
\def\QCohstack{\mathcal{QC}\!\mathit{oh}}
\def\Cohstack{\mathcal{C}\!\mathit{oh}}
\def\Spacesstack{\mathcal{S}\!\mathit{paces}}
\def\Quotfunctor{\mathrm{Quot}}
\def\Hilbfunctor{\mathrm{Hilb}}
\def\Curvesstack{\mathcal{C}\!\mathit{urves}}
\def\Polarizedstack{\mathcal{P}\!\mathit{olarized}}
\def\Complexesstack{\mathcal{C}\!\mathit{omplexes}}
% \Pic is the operator that assigns to X its picard group, usage \Pic(X)
% \Picardstack_{X/B} denotes the Picard stack of X over B
% \Picardfunctor_{X/B} denotes the Picard functor of X over B
\def\Pic{\mathop{\mathrm{Pic}}\nolimits}
\def\Picardstack{\mathcal{P}\!\mathit{ic}}
\def\Picardfunctor{\mathrm{Pic}}
\def\Deformationcategory{\mathcal{D}\!\mathit{ef}}
\newenvironment{reference}{}{}
\newenvironment{slogan}{}{}
\newenvironment{history}{}{}
\usepackage{cleveref}
\crefname{originallemma}{original lemma}{original lemmas}
\Crefname{originallemma}{Original lemma}{Original lemmas}
\makeatletter
\newlabel{spaces-properties-lemma-stalk-pullback}{{04K2}{source}{}{sourcepointer-1}{}}
\newlabel{spaces-properties-lemma-stalk-pullback@cref}{{[originallemma][0][]04K2}{[1][0][]source}{}{}{}}
\makeatother
\crefname{originallemma}{original lemma}{original lemmas}
\Crefname{originallemma}{Original lemma}{Original lemmas}
\makeatletter
\newlabel{spaces-cohomology-lemma-weak-chow}{{089J}{source}{}{sourcepointer-2}{}}
\newlabel{spaces-cohomology-lemma-weak-chow@cref}{{[originallemma][0][]089J}{[1][0][]source}{}{}{}}
\makeatother
\crefname{originallemma}{original lemma}{original lemmas}
\Crefname{originallemma}{Original lemma}{Original lemmas}
\makeatletter
\newlabel{topologies-lemma-surjective-proper-ph}{{0DES}{source}{}{sourcepointer-3}{}}
\newlabel{topologies-lemma-surjective-proper-ph@cref}{{[originallemma][0][]0DES}{[1][0][]source}{}{}{}}
\makeatother
\crefname{originallemma}{original lemma}{original lemmas}
\Crefname{originallemma}{Original lemma}{Original lemmas}
\makeatletter
\newlabel{etale-cohomology-lemma-describe-pullback-pi-ph}{{0DDW}{source}{}{sourcepointer-4}{}}
\newlabel{etale-cohomology-lemma-describe-pullback-pi-ph@cref}{{[originallemma][0][]0DDW}{[1][0][]source}{}{}{}}
\makeatother
\crefname{originallemma}{original lemma}{original lemmas}
\Crefname{originallemma}{Original lemma}{Original lemmas}
\makeatletter
\newlabel{etale-cohomology-lemma-h0-proper-over-henselian-pair}{{0A0C}{source}{}{sourcepointer-5}{}}
\newlabel{etale-cohomology-lemma-h0-proper-over-henselian-pair@cref}{{[originallemma][0][]0A0C}{[1][0][]source}{}{}{}}
\makeatother
\crefname{originallemma}{original lemma}{original lemmas}
\Crefname{originallemma}{Original lemma}{Original lemmas}
\makeatletter
\newlabel{etale-cohomology-lemma-proper-base-change-f-star}{{0A3U}{source}{}{sourcepointer-6}{}}
\newlabel{etale-cohomology-lemma-proper-base-change-f-star@cref}{{[originallemma][0][]0A3U}{[1][0][]source}{}{}{}}
\makeatother
\crefname{originalsection}{original section}{original sections}
\Crefname{originalsection}{Original section}{Original sections}
\makeatletter
\newlabel{etale-cohomology-section-proper-base-change}{{095S}{source}{}{sourcepointer-7}{}}
\newlabel{etale-cohomology-section-proper-base-change@cref}{{[originalsection][0][]095S}{[1][0][]source}{}{}{}}
\makeatother
\crefname{originallemma}{original lemma}{original lemmas}
\Crefname{originallemma}{Original lemma}{Original lemmas}
\makeatletter
\newlabel{spaces-cohomology-lemma-colimit-cohomology}{{07U6}{source}{}{sourcepointer-8}{}}
\newlabel{spaces-cohomology-lemma-colimit-cohomology@cref}{{[originallemma][0][]07U6}{[1][0][]source}{}{}{}}
\makeatother
\crefname{originallemma}{original lemma}{original lemmas}
\Crefname{originallemma}{Original lemma}{Original lemmas}
\makeatletter
\newlabel{derived-lemma-leray-acyclicity}{{015E}{source}{}{sourcepointer-9}{}}
\newlabel{derived-lemma-leray-acyclicity@cref}{{[originallemma][0][]015E}{[1][0][]source}{}{}{}}
\makeatother
\title{Proper base change for torsion etale sheaves on algebraic spaces}
\author{The Stacks Project Authors}
\date{stacks-project:spaces-more-cohomology.tex:local-snapshot-20261011; GNU FDL 1.2 or later}
\begin{document}
\maketitle
\noindent\textit{Editorial note.} Original author theorem and proof are retained. Only this target is counted; local supporting statements are not separate new entries.\par\medskip
\noindent\textit{Allowed standard prerequisites (editorial).}\begin{itemize}
\item Scheme proper base change, algebraic-space etale stalks, filtered colimits of cohomology, the acyclic-embedding criterion and Chow covers.. These are the named standard inputs; the selected main construction and its local nonroutine support remain in full.
\end{itemize}\medskip
\ifdefined\hypertarget\hypertarget{sourcepointer-1}{}\fi
\noindent\textit{Source reference.} Original lemma 04K2: focused-seeds/AGP/sources/stacks-snapshot-20261011/tags.txt, source lines 5958--5958. Stable Stacks locator for declared standard input or surrounding context. No selected local proof is replaced by this pointer.\par\medskip
\ifdefined\hypertarget\hypertarget{sourcepointer-2}{}\fi
\noindent\textit{Source reference.} Original lemma 089J: focused-seeds/AGP/sources/stacks-snapshot-20261011/tags.txt, source lines 10499--10499. Stable Stacks locator for declared standard input or surrounding context. No selected local proof is replaced by this pointer.\par\medskip
\ifdefined\hypertarget\hypertarget{sourcepointer-3}{}\fi
\noindent\textit{Source reference.} Original lemma 0DES: focused-seeds/AGP/sources/stacks-snapshot-20261011/tags.txt, source lines 16812--16812. Stable Stacks locator for declared standard input or surrounding context. No selected local proof is replaced by this pointer.\par\medskip
\ifdefined\hypertarget\hypertarget{sourcepointer-4}{}\fi
\noindent\textit{Source reference.} Original lemma 0DDW: focused-seeds/AGP/sources/stacks-snapshot-20261011/tags.txt, source lines 16781--16781. Stable Stacks locator for declared standard input or surrounding context. No selected local proof is replaced by this pointer.\par\medskip
\ifdefined\hypertarget\hypertarget{sourcepointer-5}{}\fi
\noindent\textit{Source reference.} Original lemma 0A0C: focused-seeds/AGP/sources/stacks-snapshot-20261011/tags.txt, source lines 12630--12630. Stable Stacks locator for declared standard input or surrounding context. No selected local proof is replaced by this pointer.\par\medskip
\ifdefined\hypertarget\hypertarget{sourcepointer-6}{}\fi
\noindent\textit{Source reference.} Original lemma 0A3U: focused-seeds/AGP/sources/stacks-snapshot-20261011/tags.txt, source lines 12752--12752. Stable Stacks locator for declared standard input or surrounding context. No selected local proof is replaced by this pointer.\par\medskip
\ifdefined\hypertarget\hypertarget{sourcepointer-7}{}\fi
\noindent\textit{Source reference.} Original section 095S: focused-seeds/AGP/sources/stacks-snapshot-20261011/tags.txt, source lines 11595--11595. Stable Stacks locator for declared standard input or surrounding context. No selected local proof is replaced by this pointer.\par\medskip
\ifdefined\hypertarget\hypertarget{sourcepointer-8}{}\fi
\noindent\textit{Source reference.} Original lemma 07U6: focused-seeds/AGP/sources/stacks-snapshot-20261011/tags.txt, source lines 9961--9961. Stable Stacks locator for declared standard input or surrounding context. No selected local proof is replaced by this pointer.\par\medskip
\ifdefined\hypertarget\hypertarget{sourcepointer-9}{}\fi
\noindent\textit{Source reference.} Original lemma 015E: focused-seeds/AGP/sources/stacks-snapshot-20261011/tags.txt, source lines 1505--1505. Stable Stacks locator for declared standard input or surrounding context. No selected local proof is replaced by this pointer.\par\medskip
\begin{lemma}
\label{lemma-surjective-proper}
Let $S$ be a scheme. Let $f : Y \to X$ be a surjective proper morphism
of algebraic spaces over $S$. Let $\mathcal{F}$ be a sheaf on $X_\etale$.
Then $\mathcal{F} \to f_*f^{-1}\mathcal{F}$ is injective with
image the equalizer of the two maps
$f_*f^{-1}\mathcal{F} \to g_*g^{-1}\mathcal{F}$ where
$g$ is the structure morphism $g : Y \times_X Y \to X$.
\end{lemma}

\begin{proof}
For any surjective morphism $f : Y \to X$ of algebraic spaces over $S$,
the map $\mathcal{F} \to f_*f^{-1}\mathcal{F}$ is injective.
Namely, if $\overline{x}$ is a geometric point of $X$, then we
choose a geometric point $\overline{y}$ of $Y$ lying over $\overline{x}$
and we consider
$$
\mathcal{F}_{\overline{x}} \to
(f_*f^{-1}\mathcal{F})_{\overline{x}} \to
(f^{-1}\mathcal{F})_{\overline{y}} = \mathcal{F}_{\overline{x}}
$$
See Properties of Spaces, Lemma \ref{spaces-properties-lemma-stalk-pullback}
for the last equality.

\medskip\noindent
The second statement is local on $X$ in the \'etale topology, hence we may
and do assume $Y$ is an affine scheme.

\medskip\noindent
Choose a surjective proper morphism $Z \to Y$ where $Z$ is a scheme, see
Cohomology of Spaces, Lemma \ref{spaces-cohomology-lemma-weak-chow}.
The result for $Z \to X$ implies the result for $Y \to X$.
Since $Z \to X$ is a surjective proper morphism of schemes
and hence a ph covering
(Topologies, Lemma \ref{topologies-lemma-surjective-proper-ph})
the result for $Z \to X$ follows from
\'Etale Cohomology, Lemma \ref{etale-cohomology-lemma-describe-pullback-pi-ph}
(in fact it is in some sense equivalent to this lemma).
\end{proof}

\begin{lemma}
\label{lemma-h0-proper-over-henselian-pair}
Let $(A, I)$ be a henselian pair. Let $X$ be an algebraic space over $A$
such that the structure morphism $f : X \to \Spec(A)$ is proper.
Let $i : X_0 \to X$ be the inclusion of $X \times_{\Spec(A)} \Spec(A/I)$.
For any sheaf $\mathcal{F}$ on $X_\etale$ we
have $\Gamma(X, \mathcal{F}) = \Gamma(X_0, i^{-1}\mathcal{F})$.
\end{lemma}

\begin{proof}
Choose a surjective proper morphism $Y \to X$ where $Y$ is a scheme, see
Cohomology of Spaces, Lemma \ref{spaces-cohomology-lemma-weak-chow}.
Consider the diagram
$$
\xymatrix{
\Gamma(X_0, \mathcal{F}_0) \ar[r] &
\Gamma(Y_0, \mathcal{G}_0) \ar@<1ex>[r] \ar@<-1ex>[r] &
\Gamma((Y \times_X Y)_0, \mathcal{H}_0) \\
\Gamma(X, \mathcal{F}) \ar[r] \ar[u] &
\Gamma(Y, \mathcal{G}) \ar@<1ex>[r] \ar@<-1ex>[r] \ar[u] &
\Gamma(Y \times_X Y, \mathcal{H}) \ar[u]
}
$$
Here $\mathcal{G}$, resp.\ $\mathcal{H}$ is the pullbackf or
$\mathcal{F}$ to $Y$, resp.\ $Y \times_X Y$ and the index $0$
indicates base change to $\Spec(A/I)$. By the case of schemes
(\'Etale Cohomology, Lemma
\ref{etale-cohomology-lemma-h0-proper-over-henselian-pair})
we see that the middle and right vertical arrows are bijective.
By Lemma \ref{lemma-surjective-proper} it follows that the left one is too.
\end{proof}

\begin{lemma}
\label{lemma-h0-proper-over-henselian-local}
Let $A$ be a henselian local ring. Let $X$ be an algebraic space
over $A$ such that $f : X \to \Spec(A)$
is a proper morphism. Let $X_0 \subset X$ be the fibre of
$f$ over the closed point. For any sheaf $\mathcal{F}$ on $X_\etale$ we
have $\Gamma(X, \mathcal{F}) = \Gamma(X_0, \mathcal{F}|_{X_0})$.
\end{lemma}

\begin{proof}
This is a special case of Lemma \ref{lemma-h0-proper-over-henselian-pair}.
\end{proof}

\begin{lemma}
\label{lemma-proper-base-change-f-star}
Let $S$ be a scheme. Let $f : X \to Y$ and $g : Y' \to Y$
be a morphisms of algebraic spaces over $S$. Assume $f$ is proper.
Set $X' = Y' \times_Y X$ with projections $f' : X' \to Y'$ and $g' : X' \to X$.
Let $\mathcal{F}$ be any sheaf on $X_\etale$. Then
$g^{-1}f_*\mathcal{F} = f'_*(g')^{-1}\mathcal{F}$.
\end{lemma}

\begin{proof}
The question is \'etale local on $Y'$. Choose a scheme $V$ and a surjective
\'etale morphism $V \to Y$. Choose a scheme $V'$ and a surjective \'etale
morphism $V' \to V \times_Y Y'$. Then we may replace $Y'$ by $V'$ and
$Y$ by $V$. Hence we may assume $Y$ and $Y'$ are schemes.
Then we may work Zariski locally on $Y$ and $Y'$ and hence we may
assume $Y$ and $Y'$ are affine schemes.

\medskip\noindent
Assume $Y$ and $Y'$ are affine schemes. Choose a surjective proper morphism
$h_1 : X_1 \to X$ where $X_1$ is a scheme, see
Cohomology of Spaces, Lemma \ref{spaces-cohomology-lemma-weak-chow}.
Set $X_2 = X_1 \times_X X_1$ and denote
$h_2 : X_2 \to X$ the structure morphism. Observe this is a scheme.
By the case of schemes
(\'Etale Cohomology, Lemma
\ref{etale-cohomology-lemma-proper-base-change-f-star})
we know the lemma is true for the cartesian diagrams
$$
\vcenter{
\xymatrix{
X'_1 \ar[r] \ar[d] & X_1 \ar[d] \\
Y' \ar[r] & Y
}
}
\quad\text{and}\quad
\vcenter{
\xymatrix{
X'_2 \ar[r] \ar[d] & X_2 \ar[d] \\
Y' \ar[r] & Y
}
}
$$
and the sheaves $\mathcal{F}_i = (X_i \to X)^{-1}\mathcal{F}$.
By Lemma \ref{lemma-surjective-proper} we have an exact sequence
$0 \to \mathcal{F} \to h_{1, *}\mathcal{F}_1 \to h_{2, *}\mathcal{F}_2$
and similarly for $(g')^{-1}\mathcal{F}$ because
$X'_2 = X'_1 \times_{X'} X'_1$. Hence we conlude that the
lemma is true (some details omitted).
\end{proof}

\noindent
Let $S$ be a scheme.
Let $f : Y \to X$ be a morphism of algebraic spaces over $S$. Let
$\overline{x} : \Spec(k) \to S$ be a geometric point. The fibre
of $f$ at $\overline{x}$ is the algebraic space
$Y_{\overline{x}} = \Spec(k) \times_{\overline{x}, X} Y$ over $\Spec(k)$.
If $\mathcal{F}$ is a sheaf on $Y_\etale$, then denote
$\mathcal{F}_{\overline{x}} = p^{-1}\mathcal{F}$
the pullback of $\mathcal{F}$ to $(Y_{\overline{x}})_\etale$.
Here $p : Y_{\overline{x}} \to Y$ is the projection.
In the following we will consider the set
$\Gamma(Y_{\overline{x}}, \mathcal{F}_{\overline{x}})$.

\begin{lemma}
\label{lemma-proper-pushforward-stalk}
Let $S$ be a scheme.
Let $f : Y \to X$ be a proper morphism of algebraic spaces over $S$. Let
$\overline{x} \to X$ be a geometric point.
For any sheaf $\mathcal{F}$ on $Y_\etale$
the canonical map
$$
(f_*\mathcal{F})_{\overline{x}} \longrightarrow
\Gamma(Y_{\overline{x}}, \mathcal{F}_{\overline{x}})
$$
is bijective.
\end{lemma}

\begin{proof}
This is a special case of Lemma \ref{lemma-proper-base-change-f-star}.
\end{proof}

\begin{theorem}
\label{theorem-proper-base-change}
Let $S$ be a scheme. Let
$$
\xymatrix{
X' \ar[r]_{g'} \ar[d]_{f'} & X \ar[d]^f \\
Y' \ar[r]^g & Y
}
$$
be a cartesian square of algebraic spaces over $S$.
Assume $f$ is proper.
Let $\mathcal{F}$ be an abelian torsion sheaf on $X_\etale$.
Then the base change map
$$
g^{-1}Rf_*\mathcal{F} \longrightarrow Rf'_*(g')^{-1}\mathcal{F}
$$
is an isomorphism.
\end{theorem}

\begin{proof}
This proof repeats a few of the arguments given in the proof of the
proper base change theorem for schemes. See
\'Etale Cohomology, Section \ref{etale-cohomology-section-proper-base-change}
for more details.

\medskip\noindent
The statement is \'etale local on $Y'$ and $Y$, hence we may assume
both $Y$ and $Y'$ are affine schemes. Observe that this in particular
proves the theorem in case $f$ is representable (we will use this
below).

\medskip\noindent
For every $n \geq 1$ let $\mathcal{F}[n]$ be the subsheaf of sections
of $\mathcal{F}$ annihilated by $n$. Then
$\mathcal{F} = \colim \mathcal{F}[n]$. By 
Cohomology of Spaces, Lemma \ref{spaces-cohomology-lemma-colimit-cohomology}
the functors $g^{-1}R^pf_*$ and $R^pf'_*(g')^{-1}$ commute
with filtered colimits. Hence it suffices to prove the theorem
if $\mathcal{F}$ is killed by $n$.

\medskip\noindent
Let $\mathcal{F} \to \mathcal{I}^\bullet$ be a resolution by
injective sheaves of $\mathbf{Z}/n\mathbf{Z}$-modules.
Observe that
$g^{-1}f_*\mathcal{I}^\bullet = f'_*(g')^{-1}\mathcal{I}^\bullet$
by Lemma \ref{lemma-proper-base-change-f-star}.
Applying Leray's acyclicity lemma
(Derived Categories, Lemma \ref{derived-lemma-leray-acyclicity})
we conclude it suffices to prove
$R^pf'_*(g')^{-1}\mathcal{I}^m = 0$ for $p > 0$ and $m \in \mathbf{Z}$.

\medskip\noindent
Choose a surjective proper morphism
$h : Z \to X$ where $Z$ is a scheme, see
Cohomology of Spaces, Lemma \ref{spaces-cohomology-lemma-weak-chow}.
Choose an injective map $h^{-1}\mathcal{I}^m \to \mathcal{J}$
where $\mathcal{J}$ is an injective sheaf of
$\mathbf{Z}/n\mathbf{Z}$-modules on $Z_\etale$.
Since $h$ is surjective the map $\mathcal{I}^m \to h_*\mathcal{J}$
is injective (see Lemma \ref{lemma-surjective-proper}).
Since $\mathcal{I}^m$ is injective we see that $\mathcal{I}^m$
is a direct summand of $h_*\mathcal{J}$. Thus it suffices
to prove the desired vanishing for $h_*\mathcal{J}$.

\medskip\noindent
Denote $h'$ the base change by $g$ and denote $g'' : Z' \to Z$
the projection. There is a spectral sequence
$$
E_2^{p, q} = R^pf'_* R^qh'_* (g'')^{-1}\mathcal{J}
$$
converging to $R^{p + q}(f' \circ h')_*(g'')^{-1}\mathcal{J}$.
Since $h$ and $f \circ h$ are representable (by schemes)
we know the result we want holds for them. Thus in the
spectral sequence we see that $E_2^{p, q} = 0$ for $q > 0$
and $R^{p + q}(f' \circ h')_*(g'')^{-1}\mathcal{J} = 0$
for $p + q > 0$. It follows that $E_2^{p, 0} = 0$ for $p > 0$.
Now
$$
E_2^{p, 0} = R^pf'_* h'_* (g'')^{-1}\mathcal{J} =
R^pf'_* (g')^{-1}h_*\mathcal{J}
$$
by Lemma \ref{lemma-proper-base-change-f-star}. This finishes the proof.
\end{proof}

\bibliographystyle{plain}
\bibliography{my}
\end{document}
